\appendix

\section{CIM のパラメータと時間}\label{app:timing}

表\ref{tab:timing}に，各問題の CIM の結合の強さと，\ref{sec:selection} 節の評価に使った時間を示す．G22・G55・G70 のパラメータは以前の研究で問題ごとに調整した値である（G55・G70 はポンプの立ち上げの速さも異なる）．ほかの問題は G22 の物理定数を使い，結合の強さだけを CIM 単体の平均カット値で選んだ．CIM と SA の時間は，ほかの計算がない状態で 24 スレッド（CIM は 512 本，SA は 512 本）で 2 回測った最小値を 1 本あたりに直した値である．LS・局所的な指標・内積は 1 スレッドで 512 候補を 3 回測った最小値の 1 候補あたりの値であり，評価ではこれを 24 で割った．固有ベクトルは問題ごとに 1 回だけ計算する．

\begin{table}[h]
\centering
\caption{CIM の結合の強さ $c$ と，各処理の時間．CIM は 2400 ラウンド，SA は $T_0=0.5$．}\label{tab:timing}
\small
\begin{tabular}{lrrrrrrrrr}
\toprule
 & & & CIM & \multicolumn{2}{c}{SA [ms]} & LS & 局所指標 & 内積 & 固有ベクトル \\
問題 & $n$ & $c$ & [ms] & 500 掃引 & 5000 掃引 & [$\mu$s] & [$\mu$s] & [$\mu$s] & [ms，1 回] \\
\midrule
G1 & 800 & $-0.01$ & 26.3 & 4.56 & 43.4 & 47 & 44 & 1.9 & 15 \\
G14 & 800 & $-0.05$ & 9.6 & 2.03 & 19.2 & 19 & 28 & 2.2 & 4 \\
G15 & 800 & $-0.03$ & 6.1 & 1.60 & 15.5 & 14 & 33 & 2.0 & 2 \\
G22 & 2000 & $-0.049$ & 20.5 & 4.78 & 45.0 & 42 & 96 & 4.9 & 15 \\
G23 & 2000 & $-0.05$ & 19.8 & 4.85 & 45.1 & 54 & 95 & 4.5 & 16 \\
G32 & 2000 & $-0.08$ & 10.1 & 2.74 & 27.0 & 21 & 69 & 6.5 & 11 \\
G39 & 2000 & $-0.03$ & 14.8 & 4.11 & 41.5 & 60 & 100 & 5.7 & 7 \\
G55 & 5000 & $-0.130$ & 39.8 & 9.83 & 100.2 & 111 & 352 & 16.5 & 31 \\
G70 & 10000 & $-0.153$ & 64.2 & 16.34 & 156.9 & 230 & 680 & 34.8 & 62 \\
\bottomrule
\end{tabular}
\end{table}

\section{タブー探索の高速化}\label{app:tabu}

同点を等確率に選ぶ TS では，毎反復で全頂点の利得を走査し，最大利得の非タブーの頂点から 1 つを等確率に選ぶ．重みが整数なら利得は $[-D, D]$（$D$ は重み付き次数の最大値）の整数なので，非タブーの頂点を利得の値ごとの配列（バケット）に入れ，各頂点の配列内の位置を記録しておけば，挿入と削除は定数時間でできる．最大利得のバケットは上から探し，そこから乱数 1 回で選ぶ．反転した頂点の隣接頂点だけ利得が変わるので，1 反復の費用は次数に比例する．タブーの頂点はバケットから外してリストで管理し，毎反復でリストを調べて，期間が終わった頂点をバケットに戻す．aspiration（タブーでも，反転すれば最良値を更新する手は許す）の候補も，このリストを調べるときに同時に探せばよい．最良状態のコピーは，最良状態から離れる手の直前にだけ取る．検査モードでは毎反復，バケットの内容と全走査で求めた候補の集合が一致することを確かめ，G22 と $\pm1$ 重みの疎・密なグラフで食い違いがないことを確認した．

\section{結果を見た後に変えた点}\label{app:changes}

事前登録と実行前の計画から，結果を見た後に変えた点をまとめる．いずれも変更前の出力をリポジトリに残した．
\begin{enumerate}
  \item 事前登録：SA が種を一度も超えずラベルが初期カット値と一致する条件を判定不能とした（\ref{sec:prereg} 節）．主予算の判定には影響しない．
  \item 事前登録：探索段階の「二峰性」の記述を，判定基準を当てはめた結果に合わせて訂正した（\ref{sec:prereg} 節）．
  \item G-set の検証：計算時間の都合で，結果を見る前に G77 を除いた．
  \item G-set の検証：判定不能の基準を「SA が種を一度でも超えた候補が 5\% 未満」に改めた（\ref{sec:gset} 節）．3 条件が判定不能になった．
  \item 選別の評価：費用を処理ごとに測り直し，方法ごとに使う処理だけを課した．ランダム初期化の SA を 4096 本に増やして非復元抽出にした．時間をほかの計算がない状態で測り直した（\ref{sec:selection} 節）．主条件の判定はどの版でも同じだった．
\end{enumerate}

\section{再現の手順}\label{app:repro}

すべてのスクリプトはリポジトリのルートから実行する．主な対応は次のとおりである．
\begin{itemize}\raggedright
  \item 事前登録：\texttt{docs/20261018/prereg\_eigen\_v1.md}，\texttt{scripts/benchmarks/seed\_value\_pilot.py}，\texttt{prereg\_eigen\_analysis.py}．
  \item 小規模問題：\texttt{docs/20261018/step2\_dev\_decisions.md}，\texttt{small\_landscape.py}，\texttt{small\_seed\_study.py}，\texttt{small\_seed\_analysis.py}．
  \item G-set の検証：\texttt{docs/20261018/step3\_plan.md}，\texttt{gset\_seed\_analysis.py}，\texttt{gset\_seed\_judge.py}．
  \item 選別の評価：\texttt{docs/20261018/step4\_plan.md}，\texttt{random\_init\_sa.py}，\texttt{selection\_costs.py}，\texttt{pipeline\_timing.py}，\texttt{selection\_eval.py}．
  \item GA と CIM の生成条件：\texttt{modules/GA.py}，\texttt{modules/tabu\_bucket.py}，\texttt{cim\_condition\_lineage.py}．
  \item 図：\texttt{scripts/plotting/paper\_figures.py}．
\end{itemize}
結果は \texttt{results/2026-10-18/} と \texttt{results/2026-10-19/} の下に，実験の種類と版ごとに保存した．
